\documentclass[11pt]{article}

\usepackage{amsmath, amssymb, graphicx, geometry}
\geometry{margin=1in}

\title{Environment-Responsive Coarse-Grained Bead Dynamics}
\author{VSEPR-SIM Project}
\date{\today}

\begin{document}

\maketitle

\begin{abstract}
The anisotropic surface-mapped bead, as developed in the companion document,
is a self-contained physical object: it carries a descriptor, an orientation,
and inertial properties, and it interacts with other beads through
orientation-dependent harmonic-space potentials.  However, its state is
intrinsic---it has no awareness of its local environment.  This document
introduces the \emph{environment-responsive} extension, in which each bead
acquires a set of fast observables (local density, coordination number,
orientational order) and a slow internal state variable with relaxation
dynamics and memory.  The coupling of this extended state back into the
interaction kernels enables density-dependent behaviour, spontaneous
ordering, hysteresis, and metastability.  The result is a coarse-grained
bead that does not merely exist in a simulation but \emph{participates} in
collective material behaviour.
\end{abstract}

% ============================================================================
\section{Introduction and Motivation}
% ============================================================================

The companion document \emph{Anisotropic Surface-Mapped Coarse-Grained
Beads for Molecular Systems} establishes the canonical bead state:
\begin{equation}
    B = \bigl\{
        \mathbf{R}_B,\;
        \mathbf{Q}_B,\;
        \{c_{\ell m}^{(k)}\},\;
        M_B,\;
        Q_B^{\mathrm{tot}},\;
        \mathbf{I}_B
    \bigr\}
\end{equation}
This state is sufficient for pairwise anisotropic interactions and
rigid-body dynamics.  However, it is \emph{context-blind}: a bead in vacuum
behaves identically to a bead in a dense crystal, because nothing in $B$
encodes the local surroundings.

Real molecular systems exhibit environment-dependent behaviour at every
scale:
\begin{itemize}
    \item aromatic stacking stabilises in dense phases but not in dilute gas
    \item electrostatic screening increases with local coordination
    \item packing-induced rigidity modifies effective interaction strength
    \item interfacial beads experience asymmetric environments
    \item ordering transitions emerge from local alignment reinforcement
\end{itemize}

To capture these phenomena, the bead must acquire a second layer of state
that responds to and encodes local structure.  This document defines that
layer, establishes its design constraints, and specifies its coupling to the
existing interaction machinery.

% ============================================================================
\section{Design Principles}
% ============================================================================

The environment-responsive extension is governed by four non-negotiable
rules.  Violation of any one produces a model that is clever but useless.

\subsection{Rule 1: $X_B$ Is Derived, Not Invented}

Every component of the environment state must be computable from measurable
local structure---positions, orientations, and distances of neighbouring
beads.  Arbitrary labels, ad hoc classifications, and externally imposed
phase tags are prohibited.  The bead must \emph{discover} its environment
from its neighbours, not be \emph{told} what environment it is in.

\subsection{Rule 2: Separate Fast Observables from Slow State}

The extension contains two distinct temporal classes:
\begin{itemize}
    \item \textbf{Fast observables:} recomputed every dynamics step from
          the current neighbour configuration.  These are instantaneous
          measurements, not state variables.
    \item \textbf{Slow state:} evolves via a relaxation equation with a
          finite timescale.  This is the memory layer---it integrates
          fast observations over time.
\end{itemize}
Conflating these two classes destroys the model's ability to exhibit
hysteresis and metastability.

\subsection{Rule 3: $X_B$ Does Not Replace the Descriptor}

The environment state modifies \emph{behaviour}, not \emph{identity}.
The descriptor coefficients $\{c_{\ell m}^{(k)}\}$ encode what the bead
intrinsically is.  The environment state $X_B$ encodes what the bead is
currently experiencing.  These are distinct concepts with distinct update
schedules, and the architecture must keep them separate.

\subsection{Rule 4: Keep It Low-Dimensional}

If $X_B$ becomes a 20-variable state vector, the result is a worse
atomistic model with extra steps.  The environment extension must remain
minimal: enough variables to capture the essential physics, and no more.
Every additional variable must justify its existence by producing observable
behaviour that cannot be obtained without it.

% ============================================================================
\section{Extended Bead State}
% ============================================================================

The environment-responsive bead is defined by the augmented state:
\begin{equation}
\boxed{
    \mathbb{B}_B = \bigl\{ B,\; X_B \bigr\}
}
\end{equation}
where $B$ is the canonical intrinsic state from the companion document and
\begin{equation}
    X_B = \bigl\{ \rho_B,\; C_B,\; P_{2,B},\; \eta_B \bigr\}
\end{equation}
is the environment-responsive state.

The first three components ($\rho_B$, $C_B$, $P_{2,B}$) are fast
observables---instantaneous measurements recomputed every step.  The fourth
component ($\eta_B$) is a slow internal state variable with relaxation
dynamics and memory.

\subsection{State Hierarchy}

The full bead state now spans five layers:

\begin{center}
\begin{tabular}{llll}
    \textbf{Layer} & \textbf{Contents} & \textbf{Update} & \textbf{Role} \\
    \hline
    A (Identity)      & $\{c_{\ell m}^{(k)}\}$, $M_B$, $Q_B^{\mathrm{tot}}$, $\mathbf{I}_B$
                      & Fixed unless remapped & What the bead is \\
    B (Pose)          & $\mathbf{R}_B$, $\mathbf{Q}_B$
                      & Every step & Where and how it is oriented \\
    C (Resolution)    & $\ell_{\max}$, channels, residuals
                      & At adaptation & Whether the descriptor suffices \\
    D (Provenance)    & Fragment ID, frame method
                      & At mapping only & Traceability \\
    \textbf{E (Environment)} & $\rho_B$, $C_B$, $P_{2,B}$, $\eta_B$
                      & \textbf{Every step} & \textbf{What the bead experiences} \\
\end{tabular}
\end{center}

Layer~E is the new contribution of this document.  It depends on Layer~B
(the bead needs to know where its neighbours are) but does not depend on
Layer~A (the environment state does not care about descriptor coefficients).

% ============================================================================
\section{Fast Observables}
% ============================================================================

Fast observables are instantaneous measurements of the bead's local
environment, recomputed every dynamics step from neighbour positions and
orientations.  They are not state variables---they carry no memory and
respond immediately to configuration changes.

\subsection{Local Density: $\rho_B$}

\begin{equation}
    \rho_B = \sum_{C \neq B} w(r_{BC})
\end{equation}

where $r_{BC} = |\mathbf{R}_C - \mathbf{R}_B|$ and $w(r)$ is a smooth,
normalised weighting function with compact support.

\paragraph{Physical meaning.}
$\rho_B$ measures the degree of crowding around bead $B$.  It is large in
dense packings and small in dilute environments.

\paragraph{Weighting function.}
A Gaussian with smooth cutoff provides a suitable default:
\begin{equation}
    w(r) = \begin{cases}
        \exp\!\left(-\dfrac{r^2}{2\sigma_\rho^2}\right)
        \cdot f_{\mathrm{sw}}(r;\, r_c,\, \delta)
        & r < r_c \\[4pt]
        0 & r \geq r_c
    \end{cases}
\end{equation}
where $\sigma_\rho$ controls the spatial scale of the density measurement,
$r_c$ is the cutoff radius, and $f_{\mathrm{sw}}$ is a switching function
ensuring continuity at $r_c$.

\paragraph{Gradient.}
The gradient of $\rho_B$ with respect to $\mathbf{R}_B$ is:
\begin{equation}
    \nabla_{\mathbf{R}_B} \rho_B
    = -\sum_{C \neq B} w'(r_{BC})\, \hat{r}_{BC}
\end{equation}
This is needed if $\rho_B$ couples into the potential energy, as the
resulting force contribution must be included for energy conservation.

\subsection{Coordination Number: $C_B$}

\begin{equation}
    C_B = \sum_{C \neq B} \mathbf{1}_{r_{BC} < r_c}
\end{equation}

\paragraph{Physical meaning.}
$C_B$ is the integer (or near-integer) count of neighbours within a cutoff
distance.  It provides a crude but extremely informative measure of
neighbourhood topology.

\paragraph{Smoothed variant.}
For differentiability, a soft coordination number may be used:
\begin{equation}
    C_B^{\mathrm{soft}} = \sum_{C \neq B}
    \frac{1 - (r_{BC}/r_c)^n}{1 - (r_{BC}/r_c)^m}
\end{equation}
with $m > n$ (e.g., $n = 6$, $m = 12$), following standard PLUMED-style
switching functions.

\subsection{Orientational Order Parameter: $P_{2,B}$}

\begin{equation}
    P_{2,B} = \frac{1}{N_B} \sum_{C \in \mathcal{N}_B}
    \frac{1}{2}\left(
        3\left(\hat{n}_B \cdot \hat{n}_C\right)^2 - 1
    \right)
\end{equation}

where $\hat{n}_B$ is the primary orientation axis of bead $B$ (e.g., the
normal to an aromatic plane), $\mathcal{N}_B$ is the set of neighbours
within cutoff, and $N_B = |\mathcal{N}_B|$.

\paragraph{Physical meaning.}
$P_{2,B}$ is the local second-rank Legendre order parameter.  It detects
the degree of orientational alignment between a bead and its neighbours:
\begin{itemize}
    \item $P_{2,B} \approx 1$: perfect parallel alignment (nematic order,
          $\pi$-stacking)
    \item $P_{2,B} \approx 0$: random orientational distribution (isotropic
          liquid)
    \item $P_{2,B} \approx -\tfrac{1}{2}$: perfect perpendicular alignment
          (anti-nematic)
\end{itemize}

\paragraph{Gradient.}
The gradient with respect to orientation $\mathbf{Q}_B$ generates a torque
contribution when $P_{2,B}$ couples into the potential.  Explicit evaluation
requires differentiation of $\hat{n}_B \cdot \hat{n}_C$ with respect to the
orientation frame.

\subsection{Observable Interpretation Table}

\begin{center}
\begin{tabular}{lccc}
    \textbf{Regime} & $\rho_B$ & $C_B$ & $P_{2,B}$ \\
    \hline
    Gas / dilute       & low       & low       & $\approx 0$ \\
    Isotropic liquid   & medium    & medium    & $\approx 0$ \\
    Aligned / crystal  & high      & high      & large ($>0$) \\
    Interface          & asymmetric & variable & unstable \\
\end{tabular}
\end{center}

These three observables jointly provide \emph{phase awareness} without
requiring explicit phase labels, external order parameters, or global
thermodynamic state identification.  The bead discovers its regime from
local structure alone.

% ============================================================================
\section{Slow Internal State Variable}
% ============================================================================

The fast observables respond instantaneously to configuration changes.
To capture memory, hysteresis, and metastability, the bead requires a
\emph{slow} state variable that integrates environmental information over
time.

\subsection{Definition}

The slow internal state variable $\eta_B \in [0, 1]$ evolves according to
a first-order relaxation equation:
\begin{equation}
\boxed{
    \frac{d\eta_B}{dt}
    = \frac{1}{\tau}\left(
        f(\rho_B,\, C_B,\, P_{2,B}) - \eta_B
    \right)
}
\end{equation}

where:
\begin{itemize}
    \item $\tau > 0$ is the relaxation timescale (controls memory strength)
    \item $f(\rho_B, C_B, P_{2,B}) \in [0, 1]$ is the target function
          that maps instantaneous observables to a desired internal state
    \item $\eta_B$ relaxes \emph{toward} $f$ but does not track it
          instantaneously
\end{itemize}

\subsection{Physical Interpretation}

$\eta_B$ represents a scalar measure of the bead's integrated local
environment.  Depending on the system under study, it may be interpreted as:
\begin{itemize}
    \item \textbf{stacking participation:} $\eta \approx 1$ when the bead
          has been in a well-ordered stacking environment for a sustained
          period
    \item \textbf{local ordering membership:} $\eta \approx 1$ in regions
          of persistent orientational alignment
    \item \textbf{crowding-induced rigidity:} $\eta$ increases with
          sustained high density, reflecting local stiffening
    \item \textbf{phase-like scalar:} $\eta \approx 0$ in disordered regions,
          $\eta \approx 1$ in ordered regions
\end{itemize}

The interpretation is chosen \emph{once} for a given model and held
consistently.  Attempting to encode multiple distinct physical phenomena into
a single scalar produces confusion, not generality.

\subsection{Relaxation Timescale}

The parameter $\tau$ controls the temporal response:
\begin{itemize}
    \item \textbf{$\tau \to 0$:} $\eta_B \to f$ instantaneously;
          no memory, no hysteresis
    \item \textbf{$\tau$ moderate:} $\eta_B$ lags behind fast observables;
          transient memory, damped oscillation
    \item \textbf{$\tau \to \infty$:} $\eta_B$ barely responds;
          strong memory, persistent metastability
\end{itemize}

The choice of $\tau$ is a physical modelling decision, not a numerical
parameter.  It sets the timescale over which the bead ``remembers'' its
recent environment.

\subsection{Minimal Target Function}

The simplest viable target function is a linear combination of the fast
observables:
\begin{equation}
    f(\rho_B, C_B, P_{2,B})
    = \alpha\, \hat{\rho}_B + \beta\, \hat{P}_{2,B}
\end{equation}
where $\hat{\rho}_B$ and $\hat{P}_{2,B}$ are normalised to $[0, 1]$ and
$\alpha + \beta = 1$.

With this choice, the evolution equation becomes:
\begin{equation}
    \frac{d\eta_B}{dt}
    = \frac{1}{\tau}\left(
        \alpha\, \hat{\rho}_B + \beta\, \hat{P}_{2,B} - \eta_B
    \right)
\end{equation}

This is the \textbf{minimal viable specification}: two parameters ($\alpha$,
$\beta$) control the relative importance of density and order, while $\tau$
controls memory depth.  Three parameters total.  If this does not produce
interesting behaviour, the problem is not that more variables are needed.

\subsection{Steady-State Analysis}

At steady state ($d\eta_B/dt = 0$):
\begin{equation}
    \eta_B^{\mathrm{ss}} = f(\rho_B, C_B, P_{2,B})
\end{equation}

The slow variable converges to the target function.  Out of steady state,
$\eta_B$ lags behind $f$, creating:
\begin{itemize}
    \item \textbf{hysteresis:} $\eta_B$ retains high values after the
          environment has become disordered, and vice versa
    \item \textbf{metastability:} beads in a locally ordered region resist
          disorder because their $\eta_B$ is still high, reinforcing the
          interactions that maintain order
    \item \textbf{nucleation asymmetry:} ordering is harder to initiate
          (requires sustained $f > 0$) than to maintain (high $\eta$
          stabilises itself)
\end{itemize}

\subsection{Numerical Integration}

For small $\Delta t / \tau$, the evolution equation is integrated via
forward Euler:
\begin{equation}
    \eta_B(t + \Delta t) = \eta_B(t) + \frac{\Delta t}{\tau}
    \left( f_B(t) - \eta_B(t) \right)
\end{equation}

For stiff cases ($\Delta t / \tau \gtrsim 0.5$), the implicit (exact)
solution is preferred:
\begin{equation}
    \eta_B(t + \Delta t) = f_B(t)
    + \left(\eta_B(t) - f_B(t)\right) e^{-\Delta t / \tau}
\end{equation}

Both forms are unconditionally stable for $\tau > 0$.

% ============================================================================
\section{Coupling $X_B$ into the Interaction Machinery}
% ============================================================================

The environment state is physically meaningful only if it modifies the
bead's interactions.  Three coupling strategies are defined, ordered by
increasing aggressiveness.

\subsection{Option A: Kernel Modulation (Recommended)}

The per-$\ell$ radial kernels acquire dependence on the internal state
variables of the interacting beads:
\begin{equation}
    K_k(\ell, r) \;\longrightarrow\;
    K_k(\ell, r;\, \eta_A, \eta_B)
\end{equation}

The simplest modulation is multiplicative:
\begin{equation}
    K_k(\ell, r;\, \eta_A, \eta_B)
    = K_k(\ell, r) \cdot g_k(\eta_A, \eta_B)
\end{equation}

\paragraph{Channel-specific modulations.}

\begin{itemize}
    \item \textbf{Dispersion enhancement under ordering:}
    \begin{equation}
        K_{\mathrm{disp}}(\ell, r;\, \eta_A, \eta_B)
        = K_{\mathrm{disp}}(\ell, r)
        \cdot \left(1 + \gamma_{\mathrm{disp}}\,
        \bar{\eta}_{AB}\right)
    \end{equation}
    where $\bar{\eta}_{AB} = \tfrac{1}{2}(\eta_A + \eta_B)$.
    This stabilises $\pi$-stacking when both beads are in an ordered
    environment.

    \item \textbf{Electrostatic screening under crowding:}
    \begin{equation}
        K_{\mathrm{elec}}(\ell, r;\, \eta_A, \eta_B)
        = K_{\mathrm{elec}}(\ell, r)
        \cdot \left(1 - \gamma_{\mathrm{elec}}\,
        \bar{\eta}_{AB}\right)
    \end{equation}
    This damps electrostatic anisotropy in dense, disordered regions
    where screening is expected.

    \item \textbf{Steric hardening under compression:}
    \begin{equation}
        K_{\mathrm{steric}}(\ell, r;\, \eta_A, \eta_B)
        = K_{\mathrm{steric}}(\ell, r)
        \cdot \left(1 + \gamma_{\mathrm{steric}}\,
        \bar{\eta}_{AB}\right)
    \end{equation}
    This increases effective excluded volume under sustained crowding.
\end{itemize}

\paragraph{Advantages.}
Kernel modulation is the cleanest coupling because:
\begin{itemize}
    \item it preserves the structure of the harmonic-space interaction engine
    \item it introduces no new interaction terms
    \item all existing per-$\ell$ decompositions remain valid
    \item gradients with respect to $\eta$ are straightforward
\end{itemize}

\subsection{Option B: Descriptor Scaling}

The descriptor coefficients themselves become functions of the internal
state:
\begin{equation}
    c_{\ell m}^{(k)} \;\longrightarrow\;
    c_{\ell m}^{(k)}(\eta_B)
    = c_{\ell m}^{(k)} \cdot h_k(\eta_B)
\end{equation}

Examples:
\begin{itemize}
    \item steric coefficients shrink under compression
          ($h_{\mathrm{steric}} = 1 - \delta\, \eta_B$)
    \item electrostatic anisotropy weakens in dense regions
          ($h_{\mathrm{elec}} = 1 - \delta\, \eta_B$)
\end{itemize}

\paragraph{Caution.}
Descriptor scaling is more aggressive than kernel modulation because it
modifies the bead's \emph{apparent identity} based on environment.  This
blurs the separation between identity (Layer~A) and environment (Layer~E).
Use only when kernel modulation is insufficient.

\subsection{Option C: Additive Environment Energy}

A separate energy contribution is defined that depends directly on the
environment observables:
\begin{equation}
    U_{\mathrm{env}} = \sum_B
    \Psi(\rho_B,\, C_B,\, P_{2,B},\, \eta_B)
\end{equation}

This captures:
\begin{itemize}
    \item packing penalties (high $\rho_B$ increases energy)
    \item ordering rewards (high $P_{2,B}$ decreases energy)
    \item frustration (competing density and order terms)
\end{itemize}

\paragraph{Force contribution.}
The additive environment energy contributes to forces via:
\begin{equation}
    \mathbf{F}_B^{\mathrm{env}} = -\nabla_{\mathbf{R}_B} U_{\mathrm{env}}
    = -\sum_B \frac{\partial \Psi}{\partial \rho_B}
    \nabla_{\mathbf{R}_B} \rho_B
    - \frac{\partial \Psi}{\partial P_{2,B}}
    \nabla_{\mathbf{R}_B} P_{2,B}
    - \cdots
\end{equation}

This requires gradients of all fast observables with respect to bead
positions and orientations.

\subsection{Recommended Starting Configuration}

The minimal viable coupling uses Option~A (kernel modulation) only, with a
single modulation parameter per channel:
\begin{equation}
    K_k(\ell, r;\, \eta_A, \eta_B)
    = K_k(\ell, r) \cdot \left(1 + \gamma_k\, \bar{\eta}_{AB}\right)
\end{equation}

Total new parameters: $\gamma_{\mathrm{steric}}$,
$\gamma_{\mathrm{elec}}$, $\gamma_{\mathrm{disp}}$, $\alpha$, $\beta$,
$\tau$---six parameters beyond the existing interaction model.  This is
sufficient.  If it does not produce interesting behaviour, the problem is
not that more variables are needed.

% ============================================================================
\section{Emergent Behaviour}
% ============================================================================

The coupling of environment state into the interaction kernels enables
qualitatively new phenomena that are impossible in a context-blind model.

\subsection{Spontaneous Ordering}

When $\gamma_{\mathrm{disp}} > 0$, beads that are already in an ordered
environment ($\eta \approx 1$) interact more strongly via the dispersion
channel.  This reinforces the ordering that produced high $\eta$ in the
first place, creating a positive feedback loop:

\begin{equation}
    \text{high } P_{2,B} \;\to\;
    \text{high } f_B \;\to\;
    \text{high } \eta_B \;\to\;
    \text{stronger } K_{\mathrm{disp}} \;\to\;
    \text{higher } P_{2,B}
\end{equation}

This is the mechanism for spontaneous stacking, liquid-crystal alignment,
and local crystallisation.

\subsection{Density-Dependent Behaviour}

Kernel modulation via $\rho_B$ produces:
\begin{itemize}
    \item steric hardening under compression
    \item electrostatic damping in dense environments (screening)
    \item dispersion saturation at high packing
\end{itemize}

These are physically expected behaviours that isotropic CG models achieve
only through tabulated state-dependent potentials.  Here they emerge from
the local environment coupling.

\subsection{Hysteresis and Metastability}

Because $\eta_B$ relaxes on a timescale $\tau$, the system exhibits
history-dependent behaviour:
\begin{itemize}
    \item a bead that was recently in an ordered region retains high $\eta$
          even after the ordering dissolves, temporarily maintaining
          stronger interactions
    \item a bead in a disordered region requires sustained exposure to
          order before $\eta$ rises enough to trigger reinforcement
    \item this asymmetry between ordering and disordering produces
          hysteresis in density-temperature phase behaviour
\end{itemize}

\subsection{Interface Sensitivity}

Beads at interfaces between ordered and disordered regions experience
asymmetric environments:
\begin{itemize}
    \item high $P_{2,B}$ on the ordered side, low on the disordered side
    \item intermediate $\eta_B$ due to mixed $f_B$ values
    \item the gradient in $\eta$ across the interface produces a natural
          surface tension analogue without explicit interfacial terms
\end{itemize}

\subsection{Summary of Accessible Phenomena}

\begin{center}
\begin{tabular}{lll}
    \textbf{Phenomenon} & \textbf{Mechanism} & \textbf{Key parameter} \\
    \hline
    Spontaneous stacking & $P_2 \to \eta \to K_{\mathrm{disp}}$ feedback
        & $\gamma_{\mathrm{disp}}$, $\beta$ \\
    Electrostatic screening & $\rho \to \eta \to K_{\mathrm{elec}}$ damping
        & $\gamma_{\mathrm{elec}}$, $\alpha$ \\
    Steric hardening & $\rho \to \eta \to K_{\mathrm{steric}}$ increase
        & $\gamma_{\mathrm{steric}}$, $\alpha$ \\
    Hysteresis & $\tau > 0$ delays $\eta$ response
        & $\tau$ \\
    Metastability & high $\eta$ reinforces ordering
        & $\gamma_{\mathrm{disp}}$, $\tau$ \\
    Nucleation asymmetry & ordering requires sustained $f > 0$
        & $\tau$, $\beta$ \\
    Interface sensitivity & $\nabla \eta$ across ordered/disordered boundary
        & all \\
\end{tabular}
\end{center}

% ============================================================================
\section{Formal Specification Summary}
% ============================================================================

\subsection{Extended Bead State}

\begin{equation}
    \mathbb{B}_B = \bigl\{
        \underbrace{
            \mathbf{R}_B,\;
            \mathbf{Q}_B,\;
            \{c_{\ell m}^{(k)}\},\;
            M_B,\;
            Q_B^{\mathrm{tot}},\;
            \mathbf{I}_B
        }_{B\;\text{(intrinsic)}},\;
        \underbrace{
            \rho_B,\;
            C_B,\;
            P_{2,B},\;
            \eta_B
        }_{X_B\;\text{(environment)}}
    \bigr\}
\end{equation}

\subsection{Fast Observable Definitions}

\begin{align}
    \rho_B &= \sum_{C \neq B} w(r_{BC}) \\[4pt]
    C_B &= \sum_{C \neq B} \mathbf{1}_{r_{BC} < r_c} \\[4pt]
    P_{2,B} &= \frac{1}{N_B} \sum_{C \in \mathcal{N}_B}
        \frac{1}{2}\left(3(\hat{n}_B \cdot \hat{n}_C)^2 - 1\right)
\end{align}

\subsection{Slow State Evolution}

\begin{equation}
    \frac{d\eta_B}{dt}
    = \frac{1}{\tau}\left(
        \alpha\, \hat{\rho}_B + \beta\, \hat{P}_{2,B} - \eta_B
    \right)
\end{equation}

\subsection{Kernel Modulation}

\begin{equation}
    K_k(\ell, r;\, \eta_A, \eta_B)
    = K_k(\ell, r) \cdot \left(1 + \gamma_k\, \bar{\eta}_{AB}\right)
\end{equation}

\subsection{Parameter Set}

\begin{center}
\begin{tabular}{cll}
    \textbf{Parameter} & \textbf{Role} & \textbf{Typical range} \\
    \hline
    $\alpha$ & density weight in target function & $[0, 1]$ \\
    $\beta$ & orientational order weight & $[0, 1]$,\; $\alpha + \beta = 1$ \\
    $\tau$ & relaxation timescale & $10$--$10^4$ fs \\
    $\gamma_{\mathrm{steric}}$ & steric kernel modulation & $[0, 1]$ \\
    $\gamma_{\mathrm{elec}}$ & electrostatic kernel modulation & $[-1, 0]$ \\
    $\gamma_{\mathrm{disp}}$ & dispersion kernel modulation & $[0, 2]$ \\
\end{tabular}
\end{center}

\subsection{Invariants}

\begin{itemize}
    \item $\eta_B \in [0, 1]$ at all times (clamp if target function
          overshoots)
    \item $|1 + \gamma_k \bar{\eta}| > 0$ for all channels (kernel must
          not change sign)
    \item Fast observables are recomputed before $\eta$ is updated
    \item $\eta$ is updated before interaction evaluation
    \item Energy conservation requires force contributions from
          $\partial K / \partial \eta \cdot \partial \eta / \partial \mathbf{R}$
          when $\eta$ is pose-dependent
\end{itemize}

% ============================================================================
\section{Relationship to Companion Document}
% ============================================================================

This document extends the companion
\emph{Anisotropic Surface-Mapped Coarse-Grained Beads for Molecular
Systems} without modifying it.

\begin{center}
\begin{tabular}{ll}
    \textbf{Companion document} & \textbf{This document} \\
    \hline
    Defines $B$ (intrinsic state) & Defines $X_B$ (environment state) \\
    Pairwise interaction engine & Environment-modulated kernels \\
    Static descriptor coefficients & Coefficients optionally scaled \\
    No environment awareness & Phase awareness from local observables \\
    Rigid-body dynamics only & Slow internal state with memory \\
    Deterministic mapping & Deterministic observables + relaxation \\
\end{tabular}
\end{center}

The interface is clean: this document adds Layer~E to the existing
four-layer architecture and modulates the existing interaction kernels via
multiplicative coupling.  No existing equations are changed; the extension
is purely additive.

% ============================================================================
\section{Conclusion}
% ============================================================================

The environment-responsive extension transforms the anisotropic bead from a
static descriptor carrier into a material participant.  By introducing three
fast observables and one slow internal state variable, the bead acquires
phase awareness, memory, and history-dependent interaction strength---all
from six additional parameters.

The distinction is between:
\begin{itemize}
    \item ``the model works'' --- beads interact anisotropically and move
          under force and torque
    \item ``the system behaves'' --- beads spontaneously order, resist
          disordering, screen electrostatics in dense regions, and exhibit
          hysteresis across phase boundaries
\end{itemize}

The first capability was established in the companion document.  The second
is established here.

\end{document}
